\documentclass[11pt,oneside]{article}

\usepackage[a4paper,margin=27mm,headheight=15pt]{geometry}
\usepackage[T1]{fontenc}
\usepackage[utf8]{inputenc}
\IfFileExists{libertinus.sty}{\usepackage{libertinus}}{\usepackage{lmodern}}
\usepackage{microtype}
\usepackage{amsmath,amssymb,amsthm,mathtools,mathrsfs}
\usepackage{bm}
\usepackage{booktabs,longtable,array,tabularx}
\usepackage{graphicx}
\usepackage{xcolor}
\usepackage{enumitem}
\usepackage{fancyhdr}
\usepackage{titlesec}
\usepackage{listings}
\usepackage{xurl}
\usepackage{hyperref}
\usepackage[nameinlink,noabbrev]{cleveref}

\definecolor{linkblue}{RGB}{25,75,135}
\definecolor{shadegray}{RGB}{246,247,249}
\definecolor{rulegray}{RGB}{195,201,208}
\hypersetup{
  colorlinks=true,
  linkcolor=linkblue,
  citecolor=linkblue,
  urlcolor=linkblue,
  pdftitle={Exact Geometric Tails and q-Pochhammer Extraction for Finite-Sinc Approximants},
  pdfauthor={Prepared for Vladimir Reshetnikov},
  pdfsubject={Exact dyadic extraction for Rvachev's up-function},
  pdfkeywords={Rvachev up-function, Fabius function, finite sinc product, box spline, dyadic rational, Thue-Morse, Prouhet cancellation, q-Pochhammer, q-binomial, Richardson extrapolation}
}

\pagestyle{fancy}
\fancyhf{}
\fancyhead[L]{\small Exact dyadic extraction for finite-sinc approximants}
\fancyhead[R]{\small\nouppercase{\leftmark}}
\fancyfoot[C]{\thepage}
\renewcommand{\headrulewidth}{0.4pt}

\titleformat{\section}{\Large\bfseries}{\thesection}{0.7em}{}
\titleformat{\subsection}{\large\bfseries}{\thesubsection}{0.7em}{}
\titleformat{\subsubsection}{\normalsize\bfseries}{\thesubsubsection}{0.7em}{}
\setcounter{secnumdepth}{3}
\setcounter{tocdepth}{2}
\setlist[itemize]{leftmargin=2em,itemsep=0.25em,topsep=0.4em}
\setlist[enumerate]{leftmargin=2.3em,itemsep=0.25em,topsep=0.4em}
\setlength{\emergencystretch}{2em}

\newtheorem{theorem}{Theorem}[section]
\newtheorem{proposition}[theorem]{Proposition}
\newtheorem{lemma}[theorem]{Lemma}
\newtheorem{corollary}[theorem]{Corollary}
\theoremstyle{definition}
\newtheorem{definition}[theorem]{Definition}
\newtheorem{algorithm}[theorem]{Algorithm}
\newtheorem{example}[theorem]{Example}
\theoremstyle{remark}
\newtheorem{remark}[theorem]{Remark}
\newtheorem{warning}[theorem]{Warning}

\newcommand{\R}{\mathbb R}
\newcommand{\N}{\mathbb N}
\newcommand{\Z}{\mathbb Z}
\newcommand{\Q}{\mathbb Q}
\newcommand{\F}{\mathcal F}
\newcommand{\up}{\operatorname{up}}
\newcommand{\sinc}{\operatorname{sinc}}
\newcommand{\supp}{\operatorname{supp}}
\newcommand{\Var}{\operatorname{Var}}
\newcommand{\E}{\mathsf E}
\newcommand{\dd}{\,\mathrm d}
\newcommand{\pos}[1]{\left(#1\right)_{+}}
\newcommand{\qbinom}[3]{\genfrac{[}{]}{0pt}{}{#1}{#2}_{#3}}
\newcommand{\half}{\tfrac12}
\newcommand{\tri}[1]{\binom{#1}{2}}
\newcommand{\abs}[1]{\left\lvert#1\right\rvert}
\newcommand{\norm}[1]{\left\lVert#1\right\rVert}
\newcommand{\floor}[1]{\left\lfloor#1\right\rfloor}
\newcommand{\ceil}[1]{\left\lceil#1\right\rceil}
\newcommand{\BigO}{\mathcal O}

\newenvironment{boxedremark}{%
  \par\smallskip\noindent\begin{center}\begin{minipage}{0.94\textwidth}%
  \color{black}\hrule\vspace{0.6em}\small
}{%
  \vspace{0.6em}\hrule\end{minipage}\end{center}\smallskip
}

\lstdefinestyle{pythonstyle}{
  language=Python,
  basicstyle=\ttfamily\small,
  backgroundcolor=\color{shadegray},
  frame=single,
  rulecolor=\color{rulegray},
  columns=fullflexible,
  keepspaces=true,
  showstringspaces=false,
  keywordstyle=\bfseries,
  commentstyle=\itshape,
  xleftmargin=0.7em,
  xrightmargin=0.7em,
  aboveskip=0.8em,
  belowskip=0.8em,
  breaklines=true
}

\title{\bfseries Exact Geometric Tails and \texorpdfstring{$q$}{q}-Pochhammer Extraction\\[0.3em]
\large Finite-sinc spline approximants to Rvachev's up-function at dyadic points}
\author{Prepared for Vladimir Reshetnikov\\
\small Companion to the \texttt{ProveIt/Analysis/FabiusFunction} development}
\date{2 September 2026}

\begin{document}
\maketitle

\begin{abstract}
Let
\[
 \widehat{\up_n}(t)=\prod_{k=0}^{n}\sinc\!\left(\frac{\pi t}{2^k}\right),
 \qquad \sinc z=\frac{\sin z}{z},
\]
and let \(\up\) be the inverse transform of the corresponding infinite product.
For a fixed dyadic rational \(x\in[-1,1]\), put \(q_n=\up_n(x)\).
This article proves that the convergence of \(q_n\) is eventually not merely
asymptotic: it is an exact finite sum of geometric progressions.
If \(m\) is the least nonnegative integer for which \(2^m x\in\Z\), and
\(d=\floor{m/2}\), then, for every \(n\ge m\),
\[
 q_n=\sum_{j=0}^{d}(-1)^j a_j4^{-j(n+1)}\up^{(2j)}(x)
     =\up(x)+\sum_{j=1}^{d}C_j(x)4^{-jn}.
\]
The coefficients \(a_j\in\Q_{>0}\) are the Bell--Bernoulli coefficients of
\(1/\widehat{\up}\).  The result is exact, with no hidden remainder.  A
one-stage sharpening covers \(n=m-1\) when \(m\) is odd, with the canonical
symmetric rectangle convention in the exceptional case \(m=1\).

The proof has three ingredients.  First, \(\up_n\) is a compactly supported
box spline with an exact truncated-power formula whose signs are the signed
Thue--Morse sequence.  Second, Prouhet cancellation lowers the degree of the
single spline cell centered at a dyadic point from \(n\) to at most \(m\).
Third, the omitted random tail is a scaled independent copy of the full
up-distribution.  On the finite-dimensional space of that local polynomial,
the usual all-orders deconvolution series is a polynomial in a nilpotent
differentiation operator and therefore terminates exactly.

Writing \(\rho=1/4\), the constant mode \(\up(x)\) can consequently be
recovered from only \(d+1\) consecutive rational values.  For any \(N\ge m\),
\[
 \boxed{\displaystyle
 \up(x)=\frac{1}{(\rho;\rho)_d}
 \sum_{k=0}^{d}(-1)^k\rho^{k(k+1)/2}
 \qbinom{d}{k}{\rho}\,q_{N+d-k}.}
\]
This is the first row of the inverse geometric Vandermonde system, written in
finite \(q\)-Pochhammer and Gaussian-binomial form.  Explicit formulas recover
all geometric coefficients and all even derivatives in the terminating jet.
The extraction is numerically well conditioned: the sum of the absolute
weights is less than \(1.969261\), independently of \(d\).  Exact examples and
a standard-library Python implementation accompany the article.
\end{abstract}

\tableofcontents

\section{Introduction and statement of the result}

The Fourier transform of Rvachev's normalized up-function is the dyadic sinc
product
\begin{equation}\label{eq:infinite-product}
 \Phi(t):=\widehat{\up}(t)
 =\prod_{k=0}^{\infty}\sinc\!\left(\frac{\pi t}{2^k}\right).
\end{equation}
Truncating this product after the factor with index \(n\) gives
\begin{equation}\label{eq:finite-product}
 \Phi_n(t):=\prod_{k=0}^{n}\sinc\!\left(\frac{\pi t}{2^k}\right),
 \qquad \up_n:=\F^{-1}\Phi_n.
\end{equation}
The function \(\up_n\) is a degree-\(n\) compactly supported spline.  At a
rational point its value is rational, and \(\up_n\to\up\) rapidly in every
fixed \(C^r\)-norm.

The all-orders analysis of finite dyadic sinc products naturally produces an
expansion in powers of \(4^{-n}\).  At a general point this is an asymptotic
expansion.  At a dyadic point, however, the expansion becomes a finite exact
identity.  This distinction is the central subject of the article.

\begin{definition}[Exact dyadic level]\label{def:level}
For an interior dyadic rational \(x\in(-1,1)\), its \emph{exact dyadic level}
\(m=\ell(x)\) is the least integer \(m\ge0\) such that \(2^m x\in\Z\).
Equivalently, \(x=p/2^m\) in lowest terms, with \(p\) odd when \(m>0\).
We put
\begin{equation}\label{eq:d-rho}
 d:=\floor{m/2},\qquad \rho:=\frac14.
\end{equation}
\end{definition}

The following theorem gives both the exact eventual form and its coefficients.

\begin{theorem}[Exact terminating deconvolution]\label{thm:main-differential}
Let \(x\in(-1,1)\) be dyadic of exact level \(m\), and let
\(d=\floor{m/2}\).  There exist positive rational numbers \(a_j\), independent
of \(x\) and \(n\), such that for every \(n\ge m\),
\begin{equation}\label{eq:main-differential}
 \boxed{\displaystyle
 \up_n(x)=\sum_{j=0}^{d}(-1)^j a_j4^{-j(n+1)}\up^{(2j)}(x).}
\end{equation}
They are characterized by
\begin{equation}\label{eq:a-generating-summary}
 \frac1{\Phi(z)}=\sum_{j=0}^{\infty}a_j(2\pi z)^{2j}
 \quad\text{near }z=0.
\end{equation}
In particular, \(a_0=1\), so \(\up(x)\) is the constant term.
\end{theorem}

Writing
\begin{equation}\label{eq:Cj-def}
 C_j(x):=(-1)^j a_j4^{-j}\up^{(2j)}(x),
\end{equation}
we obtain the form requested in the problem.

\begin{corollary}[Finite geometric tail]\label{cor:geometric-tail}
Under the hypotheses of \cref{thm:main-differential},
\begin{equation}\label{eq:geometric-tail}
 \boxed{\displaystyle
 q_n:=\up_n(x)=\up(x)+\sum_{j=1}^{d}C_j(x)\rho^{jn},
 \qquad n\ge m.}
\end{equation}
Thus the nonconstant terms are geometric progressions with ratios
\[
 \rho,\rho^2,\ldots,\rho^d
 =\frac14,\frac1{16},\ldots,4^{-d}.
\]
The highest coefficient \(C_d(x)\) is nonzero whenever \(d>0\); hence the
right-hand side is a polynomial of exact degree \(d\) in \(\rho^n\).
\end{corollary}

The constant mode can be read off without first computing the \(C_j\).
We use the standard notation
\begin{equation}\label{eq:q-notation-intro}
 (a;\rho)_r:=\prod_{j=0}^{r-1}(1-a\rho^j),
 \qquad
 \qbinom{r}{k}{\rho}:=
 \frac{(\rho;\rho)_r}{(\rho;\rho)_k(\rho;\rho)_{r-k}}.
\end{equation}

\begin{theorem}[Single-formula extraction]\label{thm:single-extraction}
Let \(x,m,d\) be as above and choose any integer \(N\ge m\).  Then
\begin{equation}\label{eq:single-extraction}
 \boxed{\displaystyle
 \up(x)=\frac{1}{(\rho;\rho)_d}
 \sum_{k=0}^{d}(-1)^k\rho^{k(k+1)/2}
 \qbinom{d}{k}{\rho}\,q_{N+d-k},
 \qquad \rho=\frac14.}
\end{equation}
Equivalently,
\begin{equation}\label{eq:forward-extraction}
 \up(x)=\sum_{i=0}^{d}w_{d,i}q_{N+i},
 \qquad
 w_{d,i}=(-1)^{d-i}
 \frac{\rho^{(d-i)(d-i+1)/2}}
 {(\rho;\rho)_i(\rho;\rho)_{d-i}}.
\end{equation}
At \(\rho=1/4\), if
\begin{equation}\label{eq:Ps-def}
 P_s:=\prod_{r=1}^{s}(4^r-1),\qquad P_0:=1,
\end{equation}
then the same weights have the integer-product form
\begin{equation}\label{eq:integer-weight}
 \boxed{\displaystyle
 w_{d,i}=(-1)^{d-i}\frac{4^{i(i+1)/2}}{P_iP_{d-i}}.}
\end{equation}
\end{theorem}

Formula \eqref{eq:single-extraction} is not an approximation.  When the input
values \(q_{N},\ldots,q_{N+d}\) are exact rationals, it returns the exact
rational value \(\up(x)\).

\begin{remark}[Indexing relative to the repository]
Some documents in the accompanying repository write \(p_N\) for the inverse
transform of the first \(N\) sinc factors, indexed by \(k=0,\ldots,N-1\).
The notation of the present problem has
\begin{equation}\label{eq:index-shift}
 \up_n=p_{n+1}.
\end{equation}
Keeping this one-step shift explicit prevents the common confusion between
\(4^{-n}\) and \(4^{-(n+1)}\) in the differential formula.
\end{remark}

\section{Fourier, convolution, and probability normalizations}

We use the Fourier convention
\begin{equation}\label{eq:fourier-convention}
 \widehat f(t)=\int_{\R}f(x)e^{-2\pi ixt}\dd x,
 \qquad
 f(x)=\int_{\R}\widehat f(t)e^{2\pi ixt}\dd t
\end{equation}
whenever the inversion formula applies classically, and otherwise in the
usual distributional sense.

For \(k\ge0\), define the centered box density
\begin{equation}\label{eq:box-density}
 b_k(x):=2^k\mathbf 1_{[-2^{-k-1},\,2^{-k-1}]}(x).
\end{equation}
Its total mass is one and
\begin{equation}\label{eq:box-transform}
 \widehat b_k(t)
 =\frac{\sin(\pi t/2^k)}{\pi t/2^k}
 =\sinc\!\left(\frac{\pi t}{2^k}\right).
\end{equation}
Consequently,
\begin{equation}\label{eq:finite-convolution}
 \up_n=b_0*b_1*\cdots*b_n.
\end{equation}
This proves at once that \(\up_n\) is nonnegative, has integral one, and is a
piecewise polynomial of degree \(n\).  Its support radius is
\begin{equation}\label{eq:An}
 A_n:=\sum_{k=0}^{n}2^{-k-1}=1-2^{-n-1},
 \qquad \supp\up_n=[-A_n,A_n].
\end{equation}

There is an equivalent probabilistic model.  Let \(U_0,U_1,\ldots\) be
independent random variables, each uniform on \([-1,1]\), and set
\begin{equation}\label{eq:random-sums}
 S_n:=\sum_{k=0}^{n}2^{-k-1}U_k,
 \qquad
 S:=\sum_{k=0}^{\infty}2^{-k-1}U_k.
\end{equation}
Then \(\up_n\) is the density of \(S_n\), while \(\up\) is the density of
\(S\).  Splitting the series after the \(n\)-th term gives the exact
independent-copy identity
\begin{equation}\label{eq:head-tail-random}
 S\ \stackrel{d}=\ S_n+\delta_nS',
 \qquad \delta_n:=2^{-n-1},
\end{equation}
where \(S'\) is independent of \(S_n\) and has the same distribution as
\(S\).  If
\begin{equation}\label{eq:dilation}
 (D_\delta f)(x):=\frac1\delta f\!\left(\frac{x}{\delta}\right),
\end{equation}
then \eqref{eq:head-tail-random} is the convolution identity
\begin{equation}\label{eq:tail-convolution}
 \boxed{\displaystyle \up=\up_n*D_{\delta_n}\up.}
\end{equation}
On the Fourier side this is
\begin{equation}\label{eq:tail-fourier}
 \Phi(t)=\Phi_n(t)\Phi(\delta_nt),
 \qquad
 \Phi_n(t)=\frac{\Phi(t)}{\Phi(\delta_nt)}.
\end{equation}
The second quotient is the source of the deconvolution expansion.  The key
point of this article is that, at a dyadic evaluation point, the corresponding
operator series acts on a finite-dimensional polynomial space and becomes an
exact finite sum.

\section{The exact finite-spline formula}

\subsection{Truncated powers and Thue--Morse signs}

Let
\begin{equation}\label{eq:tau-def}
 \tau_k:=(-1)^{s_2(k)},
\end{equation}
where \(s_2(k)\) is the number of ones in the binary expansion of \(k\).
Thus
\[
 (\tau_k)_{k\ge0}=1,-1,-1,1,-1,1,1,-1,\ldots.
\]
We use \(y_+^r=y^r\) for \(y>0\) and \(0\) for \(y\le0\).  For \(r=0\),
this convention gives the right-continuous Heaviside representative; only
\(\up_0\) depends on the endpoint convention.

\begin{proposition}[Exact truncated-power representation]\label{prop:truncated-power}
For every \(n\ge0\),
\begin{equation}\label{eq:truncated-power}
 \boxed{\displaystyle
 \up_n(x)=\frac{2^{n(n+1)/2}}{n!}
 \sum_{k=0}^{2^{n+1}-1}\tau_k
 \pos{x+A_n-\frac{k}{2^n}}^{n},
 \qquad A_n=1-2^{-n-1}.}
\end{equation}
The knots therefore lie on the uniform grid
\begin{equation}\label{eq:knot-grid}
 -A_n+\frac{k}{2^n},\qquad 0\le k\le2^{n+1}-1.
\end{equation}
\end{proposition}

\begin{proof}
Write each box as a normalized difference of two translated Heaviside
functions.  Convolving \(n+1\) such differences and using the elementary
identity that the \((n+1)\)-fold integral of a point mass is
\(x_+^n/n!\), one obtains
\begin{equation}\label{eq:subset-form}
 \up_n(x)=\frac{\prod_{j=0}^{n}2^j}{n!}
 \sum_{\varepsilon\in\{0,1\}^{n+1}}
 (-1)^{\varepsilon_0+\cdots+\varepsilon_n}
 \pos{x+A_n-\sum_{j=0}^{n}\varepsilon_j2^{-j}}^n.
\end{equation}
Now put
\[
 k:=\sum_{j=0}^{n}\varepsilon_j2^{n-j}.
\]
As \(\varepsilon\) ranges over the binary cube, \(k\) ranges once over
\(0,\ldots,2^{n+1}-1\),
\[
 \sum_{j=0}^{n}\varepsilon_j2^{-j}=\frac{k}{2^n},
 \qquad
 (-1)^{\sum_j\varepsilon_j}=\tau_k,
\]
and \(\prod_{j=0}^{n}2^j=2^{n(n+1)/2}\).  Substitution gives
\eqref{eq:truncated-power}.
\end{proof}

The scaled coordinate
\begin{equation}\label{eq:y-scaled}
 y:=2^n(x+A_n)
\end{equation}
places consecutive knots one unit apart.  Since only terms with \(k<y\) are
active, \cref{prop:truncated-power} is equivalently
\begin{equation}\label{eq:cell-formula}
 \boxed{\displaystyle
 \up_n(x)=\frac{2^{-n(n-1)/2}}{n!}
 \sum_{0\le k\le\floor y}\tau_k(y-k)^n,}
\end{equation}
away from knots.  Formula \eqref{eq:cell-formula} makes both exact rationality
and the local cancellation mechanism transparent.

\begin{corollary}[Rational finite samples]\label{cor:rational-samples}
If \(x\in\Q\), then \(\up_n(x)\in\Q\) for every \(n\ge0\), subject only to
the chosen representative at the two jumps of \(\up_0\).
\end{corollary}

\subsection{Prouhet cancellation}

The signed Thue--Morse sequence annihilates low-degree polynomials on every
complete dyadic block.

\begin{lemma}[Prouhet--Thue--Morse cancellation]\label{lem:prouhet}
For every integer \(r\ge1\) and every polynomial \(P\) of degree less than
\(r\),
\begin{equation}\label{eq:prouhet-polynomial}
 \sum_{j=0}^{2^r-1}\tau_jP(j)=0.
\end{equation}
Equivalently,
\begin{equation}\label{eq:prouhet-moments}
 \sum_{j=0}^{2^r-1}\tau_jj^s=0,
 \qquad 0\le s<r.
\end{equation}
\end{lemma}

\begin{proof}
The finite generating polynomial factors as
\begin{equation}\label{eq:TM-factor}
 \sum_{j=0}^{2^r-1}\tau_jz^j
 =\prod_{h=0}^{r-1}(1-z^{2^h}).
\end{equation}
The right-hand side has a zero of order \(r\) at \(z=1\).  Applying
\((z\frac{\dd}{\dd z})^s\) and then setting \(z=1\) gives
\eqref{eq:prouhet-moments} for \(s<r\).  Linearity gives
\eqref{eq:prouhet-polynomial}.
\end{proof}

We will also use the no-carry multiplicativity
\begin{equation}\label{eq:no-carry}
 \tau_{2^r\ell+j}=\tau_\ell\tau_j,
 \qquad 0\le j<2^r.
\end{equation}

\section{Dyadic cell centers and local degree collapse}

Let \(x\in(-1,1)\) have exact level \(m\), and define the odd integer
\begin{equation}\label{eq:A-dyadic}
 A:=2^m(x+1).
\end{equation}
For \(m>0\), writing \(x=p/2^m\) in lowest terms gives \(A=p+2^m\), which
is odd.  The same is true for \(m=0,x=0\), where \(A=1\).

Fix \(n\ge m\) and put \(r=n-m\).  At the point \(x\), the scaled cell
coordinate \eqref{eq:y-scaled} is
\begin{equation}\label{eq:midpoint-coordinate}
 y=2^n(x+A_n)
   =2^n(x+1)-\frac12
   =A2^r-\frac12.
\end{equation}
Thus \(x\) is exactly the midpoint between two consecutive knots of
\(\up_n\).  The cell has half-width
\begin{equation}\label{eq:cell-halfwidth}
 \delta_n=2^{-n-1},
\end{equation}
which is also the scale of the omitted random tail in
\eqref{eq:head-tail-random}.

\begin{theorem}[Local degree collapse]\label{thm:degree-collapse}
Let \(x\) have exact dyadic level \(m\).  For every \(n\ge m\), the
restriction of \(\up_n\) to the open cell
\begin{equation}\label{eq:center-cell}
 (x-\delta_n,x+\delta_n)
\end{equation}
is a polynomial of degree at most \(m\), even though \(\up_n\) has global
spline degree \(n\).
\end{theorem}

\begin{proof}
Write \(h\) for the local coordinate, \(\abs h<\delta_n\).  Then
\(\floor{2^n(x+h+A_n)}=A2^r-1\), and \eqref{eq:cell-formula} gives, up to the
nonzero constant \(2^{-n(n-1)/2}/n!\),
\begin{equation}\label{eq:block-before}
 \sum_{k=0}^{A2^r-1}\tau_k
 \left(A2^r-\frac12+2^nh-k\right)^n.
\end{equation}
Split \(k=2^r\ell+j\), where \(0\le\ell<A\) and \(0\le j<2^r\).  By
\eqref{eq:no-carry}, the expression becomes
\begin{equation}\label{eq:block-after}
 \sum_{\ell=0}^{A-1}\tau_\ell
 \sum_{j=0}^{2^r-1}\tau_j
 \left(2^r(A-\ell)-\frac12-j+2^nh\right)^n.
\end{equation}
The coefficient of \(h^s\) is proportional to
\begin{equation}\label{eq:h-coefficient}
 \sum_{\ell=0}^{A-1}\tau_\ell
 \sum_{j=0}^{2^r-1}\tau_j
 \left(2^r(A-\ell)-\frac12-j\right)^{n-s}.
\end{equation}
If \(s>m\), then
\[
 n-s<n-m=r.
\]
For each fixed \(\ell\), the inner summand in \eqref{eq:h-coefficient} is a
polynomial in \(j\) of degree less than \(r\), so it vanishes by
\cref{lem:prouhet}.  Every coefficient of \(h^s\) with \(s>m\) is therefore
zero.
\end{proof}

\begin{remark}[What the cancellation really does]
The global degree grows with \(n\), but a fixed dyadic point organizes the
active truncated powers into complete blocks of length \(2^{n-m}\).  Each
extra sinc factor creates one extra binary cancellation moment.  The increase
in global polynomial degree is exactly offset by the increase in Prouhet
order, leaving a local degree bounded by the fixed denominator exponent
\(m\).
\end{remark}

\section{Bell--Bernoulli deconvolution coefficients}

\subsection{The moment-generating function}

Let
\begin{equation}\label{eq:mgf}
 M(s):=\int_{-1}^{1}e^{sy}\up(y)\dd y
      =\Phi\!\left(\frac{is}{2\pi}\right).
\end{equation}
Since \(\up\) is even, \(M\) is an even entire function and
\begin{equation}\label{eq:mgf-moments}
 M(s)=\sum_{r=0}^{\infty}\frac{\mu_{2r}}{(2r)!}s^{2r},
 \qquad
 \mu_{2r}:=\int_{-1}^{1}y^{2r}\up(y)\dd y.
\end{equation}

The classical Bernoulli expansion
\begin{equation}\label{eq:log-sinc-recip}
 \log\frac{z}{\sin z}
 =\sum_{r=1}^{\infty}
 \frac{2^{2r-1}\abs{B_{2r}}}{r(2r)!}z^{2r}
\end{equation}
combined with the geometric sum over the dyadic scales yields
\begin{equation}\label{eq:alpha-expansion}
 \log\frac1{\Phi(z)}
 =\sum_{r=1}^{\infty}\alpha_r(2\pi z)^{2r},
 \qquad
 \boxed{\displaystyle
 \alpha_r=\frac{\abs{B_{2r}}}
 {2r(2r)!\left(1-2^{-2r}\right)}.}
\end{equation}
Define \(a_j\) by
\begin{equation}\label{eq:a-def}
 \frac1{\Phi(z)}
 =\exp\!\left(\sum_{r=1}^{\infty}\alpha_r(2\pi z)^{2r}\right)
 =\sum_{j=0}^{\infty}a_j(2\pi z)^{2j}.
\end{equation}
Then every \(a_j\) is a positive rational number.  In terms of complete
exponential Bell polynomials,
\begin{equation}\label{eq:a-bell}
 \boxed{\displaystyle
 a_j=\frac1{j!}Y_j\bigl(1!\alpha_1,2!\alpha_2,\ldots,j!\alpha_j\bigr).}
\end{equation}
Equivalently,
\begin{equation}\label{eq:a-partition}
 a_j=\sum_{k_1+2k_2+\cdots+jk_j=j}
 \prod_{r=1}^{j}\frac{\alpha_r^{k_r}}{k_r!}.
\end{equation}
The most convenient exact recurrence is
\begin{equation}\label{eq:a-recurrence}
 \boxed{\displaystyle
 a_0=1,
 \qquad
 ja_j=\sum_{r=1}^{j}r\alpha_r a_{j-r}\quad(j\ge1).}
\end{equation}

The first coefficients are
\begin{equation}\label{eq:first-a}
 \begin{aligned}
 a_0&=1,
 &a_1&=\frac1{18},
 &a_2&=\frac{31}{16200},\\
 a_3&=\frac{2347}{42865200},
 &a_4&=\frac{1904369}{1311675120000},
 &a_5&=\frac{3309760579}{88561680752160000}.
 \end{aligned}
\end{equation}
Substitution of \(z=is/(2\pi)\) into \eqref{eq:a-def} gives the reciprocal
moment series
\begin{equation}\label{eq:reciprocal-M}
 \boxed{\displaystyle
 \frac1{M(s)}=\sum_{j=0}^{\infty}(-1)^ja_js^{2j}.}
\end{equation}

\subsection{An exact inverse on polynomials}

Let \(D=\frac{\dd}{\dd z}\).  If \(P\) is a polynomial, define
\begin{equation}\label{eq:Q-polynomial}
 Q(z):=\int_{-1}^{1}P(z-\delta y)\up(y)\dd y.
\end{equation}
Taylor expansion terminates, and evenness of \(\up\) gives the exact operator
identity
\begin{equation}\label{eq:operator-M}
 Q=M(\delta D)P.
\end{equation}

\begin{lemma}[Finite polynomial deconvolution]\label{lem:polynomial-inverse}
If \(P\) has degree at most \(m\), and \(Q=M(\delta D)P\), then
\begin{equation}\label{eq:polynomial-inverse}
 \boxed{\displaystyle
 P=\sum_{j=0}^{\floor{m/2}}(-1)^ja_j\delta^{2j}Q^{(2j)}.}
\end{equation}
More generally, for \(0\le r\le m\),
\begin{equation}\label{eq:polynomial-inverse-derivative}
 P^{(r)}(0)=
 \sum_{j=0}^{\floor{(m-r)/2}}(-1)^ja_j\delta^{2j}Q^{(r+2j)}(0).
\end{equation}
\end{lemma}

\begin{proof}
On the vector space of polynomials of degree at most \(m\), the operator
\(D\) is nilpotent: \(D^{m+1}=0\).  Therefore every formal power series in
\(D\) is actually a finite polynomial in \(D\), and multiplication of formal
series is legitimate without any convergence question.  Equations
\eqref{eq:operator-M} and \eqref{eq:reciprocal-M} give
\[
 P=M(\delta D)^{-1}Q
  =\sum_{j\ge0}(-1)^ja_j\delta^{2j}D^{2j}Q.
\]
All terms with \(2j>m\) vanish, proving \eqref{eq:polynomial-inverse}.
Differentiating \(r\) times and evaluating at zero gives
\eqref{eq:polynomial-inverse-derivative}.
\end{proof}

\begin{boxedremark}
The exactness in \cref{lem:polynomial-inverse} is stronger than observing that
all sufficiently high coefficients of an asymptotic expansion vanish at a
dyadic point.  A vanishing formal tail alone would not rule out a remainder
that is smaller than every power of \(4^{-n}\).  Here there is no such issue:
the inverse series is evaluated on a nilpotent operator, so it is literally a
finite algebraic identity.
\end{boxedremark}

\section{Proof of exact termination at dyadic points}

\subsection{The centered-cell case}

Fix a dyadic \(x\) of level \(m\) and a stage \(n\ge m\).  By
\cref{thm:degree-collapse}, there is a polynomial \(P\) of degree at most
\(m\) such that
\begin{equation}\label{eq:P-local}
 P(z)=\up_n(x+z),\qquad \abs z<\delta_n.
\end{equation}
Define \(Q\) from \(P\) by \eqref{eq:Q-polynomial} with
\(\delta=\delta_n\).

For \(0\le r\le m\), the convolution identity
\eqref{eq:tail-convolution}, interpreted distributionally when necessary,
gives
\begin{equation}\label{eq:derivative-convolution}
 \up^{(r)}(x)=
 \int_{-1}^{1}\up_n^{(r)}(x-\delta_ny)\up(y)\dd y.
\end{equation}
For \(-1<y<1\), the point \(x-\delta_ny\) lies strictly inside the centered
cell, so \(\up_n^{(r)}=P^{(r)}\) there.  The two endpoints have measure zero.
Consequently,
\begin{equation}\label{eq:Q-jet}
 Q^{(r)}(0)=\up^{(r)}(x),\qquad 0\le r\le m.
\end{equation}
Applying \cref{lem:polynomial-inverse} at \(z=0\), and using
\(P(0)=\up_n(x)\), gives
\[
 \up_n(x)=\sum_{j=0}^{\floor{m/2}}
 (-1)^ja_j\delta_n^{2j}\up^{(2j)}(x).
\]
Since \(\delta_n^{2j}=4^{-j(n+1)}\), this proves
\cref{thm:main-differential}.

The same proof yields a derivative version.

\begin{corollary}[Terminating laws for finite-spline derivatives]\label{cor:derivative-tail}
Let \(0\le r\le m\).  For every \(n\ge m\),
\begin{equation}\label{eq:derivative-tail}
 \boxed{\displaystyle
 \up_n^{(r)}(x)=
 \sum_{j=0}^{\floor{(m-r)/2}}(-1)^ja_j
 4^{-j(n+1)}\up^{(r+2j)}(x).}
\end{equation}
Thus the derivative sequence also has a finite exact geometric tail, with at
most \(\floor{(m-r)/2}\) nonconstant modes.
\end{corollary}

\subsection{A one-stage sharpening at odd dyadic levels}

The universal centered-cell proof starts at \(n=m\).  When \(m\) is odd, the
same identity already holds one stage earlier.

\begin{proposition}[Odd-level onset]\label{prop:odd-onset}
Let \(m=2d+1\ge3\), and let \(x\) have exact level \(m\).  Then
\eqref{eq:main-differential} also holds at \(n=m-1=2d\).  Hence the geometric
law \eqref{eq:geometric-tail} is valid for every
\begin{equation}\label{eq:nu-onset}
 n\ge\nu(x):=2\floor{m/2}.
\end{equation}
For \(m=1\), the same statement at \(n=0\) holds if the inverse transform of
the first sinc factor is assigned its canonical symmetric value \(1/2\) at
\(x=\pm1/2\).  Avoiding this endpoint convention, one may simply start at
\(n=m=1\).
\end{proposition}

\begin{proof}
Put \(n=m-1=2d\).  In this case \(x\) is a knot of the degree-\(n\) spline
\(\up_n\), and the interval \([x-\delta_n,x+\delta_n]\) lies in the union of the two
adjacent knot cells, meeting each in a half-cell.  Let \(P_-\) and \(P_+\) be the polynomial pieces in the
local coordinate \(z\), on the left and right respectively.  Since a
convolution of \(n+1\) boxes is \(C^{n-1}\), the two pieces have the form
\begin{equation}\label{eq:two-pieces}
 P_-(z)=R(z)+a_-z^n,
 \qquad
 P_+(z)=R(z)+a_+z^n,
 \qquad \deg R\le n-1.
\end{equation}
Define their top-coefficient average
\begin{equation}\label{eq:P-star}
 P_*(z):=R(z)+\frac{a_-+a_+}{2}z^n.
\end{equation}
Because \(n\) is even, for every even \(r\le n\), evenness of \(\up\) shows
that the contributions from the two half-intervals satisfy
\begin{align}\label{eq:averaged-jet}
 \up^{(r)}(x)
 &=\int_0^1
 \left(P_-^{(r)}(-\delta_ny)+P_+^{(r)}(\delta_ny)\right)\up(y)\dd y\\
 &=\int_{-1}^{1}P_*^{(r)}(-\delta_ny)\up(y)\dd y.
\end{align}
Moreover, continuity gives \(P_*(0)=P_-(0)=P_+(0)=\up_n(x)\).  Apply the
finite polynomial inverse \cref{lem:polynomial-inverse} to \(P_*\), whose
degree is \(n=2d\).  Only even derivatives occur, and
\eqref{eq:averaged-jet} supplies exactly the required jet of \(\up\).  The
result is \eqref{eq:main-differential} at \(n=m-1\).

For \(m=1,n=0\), the same argument averages the two constant pieces of a
rectangle.  Symmetric Fourier inversion assigns their half-sum at the jump.
\end{proof}

\subsection{Vanishing derivatives and the highest geometric mode}

Repeated differentiation of Rvachev's functional-differential equation
\begin{equation}\label{eq:up-differential}
 \up'(x)=2\bigl(\up(2x+1)-\up(2x-1)\bigr)
\end{equation}
gives an exact Thue--Morse formula for every derivative.

\begin{lemma}[Derivative dilation formula]\label{lem:derivative-formula}
For every integer \(r\ge0\),
\begin{equation}\label{eq:derivative-formula}
 \boxed{\displaystyle
 \up^{(r)}(x)=2^{r(r+1)/2}
 \sum_{k=0}^{2^r-1}\tau_k\,
 \up\!\left(2^rx+2^r-1-2k\right).}
\end{equation}
\end{lemma}

\begin{proof}
The case \(r=0\) is tautological, and the case \(r=1\) is
\eqref{eq:up-differential}.  Differentiate the formula for order \(r\), apply
\eqref{eq:up-differential} to each term, and split each old index into the two
new binary indices \(k\) and \(k+2^r\).  The identities
\(\tau_{k+2^r}=-\tau_k\) and
\(2^{r(r+1)/2}2^r2=2^{(r+1)(r+2)/2}\) give the order-\(r+1\) formula.
\end{proof}

The values needed below follow directly from the already proved terminating
formula.  At \(x=0\), level \(m=0\) gives \(\up(0)=\up_0(0)=1\).  At
\(x=\pm1/2\), level \(m=1\) gives \(\up(\pm1/2)=\up_1(\pm1/2)=1/2\), the
last value being immediate from the convolution of the boxes of half-widths
\(1/2\) and \(1/4\).

\begin{proposition}[Dyadic derivative cutoff and nonzero top even jet]\label{prop:top-jet}
Let \(x=p/2^m\in(-1,1)\) have exact level \(m\), and put
\(A=2^m(x+1)\).  Then
\begin{equation}\label{eq:derivative-cutoff}
 \up^{(r)}(x)=0\qquad(r>m).
\end{equation}
If \(m=2d\) is even and positive, then
\begin{equation}\label{eq:top-even-evenlevel}
 \up^{(m)}(x)=2^{m(m+1)/2}\tau_{(A-1)/2}\ne0.
\end{equation}
If \(m=2d+1\) is odd, then
\begin{equation}\label{eq:top-even-oddlevel}
 \up^{(m-1)}(x)=
 2^{m(m-1)/2-1}\tau_{\floor{(A-1)/4}}\ne0.
\end{equation}
\end{proposition}

\begin{proof}
For \(r>m\), the number \(2^rx\) is an even integer.  Every argument of
\(\up\) in \eqref{eq:derivative-formula} is then an odd integer.  The only
odd integers in \([-1,1]\) are \(\pm1\), and \(\up(\pm1)=0\); hence the
whole sum vanishes.

For \(r=m\), the arguments in \eqref{eq:derivative-formula} are even
integers.  Exactly one is zero, at \(k=(A-1)/2\), and every other argument
lies outside \((-1,1)\).  Since \(\up(0)=1\), this proves
\eqref{eq:top-even-evenlevel} when \(m\) is even.

If \(m\) is odd, take \(r=m-1\).  The arguments are half-integers separated
by two.  Exactly one equals \(1/2\) or \(-1/2\), at
\(k=\floor{(A-1)/4}\).  Since \(\up(\pm1/2)=1/2\), formula
\eqref{eq:top-even-oddlevel} follows.
\end{proof}

Since \(a_d>0\), \cref{prop:top-jet} proves the nonvanishing assertion in
\cref{cor:geometric-tail}.  It also gives an explicit sign and magnitude for
the slowest-to-appear, fastest-decaying top mode \(C_d(x)\).

\section{The geometric Vandermonde system}

Assume from now on that \(N\ge m\), so no convention at \(\up_0\) is
involved.  Put
\begin{equation}\label{eq:Bj}
 B_j:=C_j(x)\rho^{jN},\qquad B_0=\up(x).
\end{equation}
The \(d+1\) consecutive samples satisfy
\begin{equation}\label{eq:vandermonde-equations}
 q_{N+i}=\sum_{j=0}^{d}B_j\rho^{ij},
 \qquad 0\le i\le d.
\end{equation}
In matrix form,
\begin{equation}\label{eq:vandermonde-matrix}
 \begin{pmatrix}
 q_N\\q_{N+1}\\\vdots\\q_{N+d}
 \end{pmatrix}
 =
 \begin{pmatrix}
 1&1&1&\cdots&1\\
 1&\rho&\rho^2&\cdots&\rho^d\\
 1&\rho^2&\rho^4&\cdots&\rho^{2d}\\
 \vdots&\vdots&\vdots&&\vdots\\
 1&\rho^d&\rho^{2d}&\cdots&\rho^{d^2}
 \end{pmatrix}
 \begin{pmatrix}
 B_0\\B_1\\\vdots\\B_d
 \end{pmatrix}.
\end{equation}
The determinant is
\begin{equation}\label{eq:vandermonde-det}
 \prod_{0\le j<k\le d}(\rho^k-\rho^j)\ne0,
\end{equation}
so the coefficients are uniquely determined.

\subsection{The first inverse row by Lagrange interpolation}

Define the polynomial
\begin{equation}\label{eq:f-polynomial}
 f(z):=\sum_{j=0}^{d}B_jz^j.
\end{equation}
Then \(f(\rho^i)=q_{N+i}\), while \(f(0)=B_0=\up(x)\).  Lagrange
interpolation at the geometric nodes \(1,\rho,\ldots,\rho^d\) therefore gives
\begin{equation}\label{eq:lagrange-zero}
 \up(x)=\sum_{i=0}^{d}
 \left(\prod_{\substack{0\le\ell\le d\\\ell\ne i}}
 \frac{-\rho^\ell}{\rho^i-\rho^\ell}\right)q_{N+i}.
\end{equation}
For \(\ell<i\), the corresponding product factors are
\(1/(1-\rho^{i-\ell})\); for \(\ell>i\), they are
\(-\rho^{\ell-i}/(1-\rho^{\ell-i})\).  Their product is exactly
\eqref{eq:forward-extraction}.  Replacing \(i\) by \(d-k\) and using the
Gaussian-binomial definition \eqref{eq:q-notation-intro} proves
\eqref{eq:single-extraction}.

The operator form of the same argument is particularly compact.  Let the
forward shift act by
\begin{equation}\label{eq:shift}
 (\E q)_n:=q_{n+1}.
\end{equation}
Each factor \(\E-\rho^j\) kills the geometric mode \(\rho^{jn}\).  Hence
\begin{equation}\label{eq:constant-projector}
 \boxed{\displaystyle
 \up(x)=\frac{1}{(\rho;\rho)_d}
 \left[\prod_{j=1}^{d}(\E-\rho^j)q\right]_N.}
\end{equation}
Expanding
\begin{equation}\label{eq:q-product-expand}
 \prod_{j=1}^{d}(\E-\rho^j)
 =\E^d(\rho\E^{-1};\rho)_d
\end{equation}
with the finite \(q\)-binomial theorem produces
\eqref{eq:single-extraction} in one line.

\subsection{First extraction formulas}

For small \(d\), the formulas are
\begin{align}
 d=0:\quad&\up(x)=q_N,\label{eq:weights-d0}\\
 d=1:\quad&\up(x)=\frac{-q_N+4q_{N+1}}{3},\label{eq:weights-d1}\\
 d=2:\quad&\up(x)=\frac{q_N-20q_{N+1}+64q_{N+2}}{45},\label{eq:weights-d2}\\
 d=3:\quad&\up(x)=
 \frac{-q_N+84q_{N+1}-1344q_{N+2}+4096q_{N+3}}{2835},\label{eq:weights-d3}\\
 d=4:\quad&\up(x)=\frac{q_N-340q_{N+1}+22848q_{N+2}
 -348160q_{N+3}+1048576q_{N+4}}{722925}.
 \label{eq:weights-d4}
\end{align}
The required number of samples grows only every second time the dyadic
denominator exponent grows:
\begin{center}
\begin{tabular}{c c c c}
\toprule
exact level \(m\)&number of modes \(d\)&samples needed&guaranteed start\\
\midrule
0 or 1&0&1&\(N\ge m\)\\
2 or 3&1&2&\(N\ge m\)\\
4 or 5&2&3&\(N\ge m\)\\
6 or 7&3&4&\(N\ge m\)\\
8 or 9&4&5&\(N\ge m\)\\
\bottomrule
\end{tabular}
\end{center}
For odd \(m\), \cref{prop:odd-onset} permits the first window to start at
\(N=m-1\), apart from the harmless \(m=1\) point-value convention.

\subsection{Uniform stability of the extraction}

Although the weights alternate in sign, their total variation stays bounded
as \(d\to\infty\).

\begin{proposition}[Lebesgue factor of the geometric extraction]\label{prop:stability}
For \(0<\rho<1\),
\begin{equation}\label{eq:stability-factor}
 \boxed{\displaystyle
 \Lambda_d(\rho):=\sum_{i=0}^{d}\abs{w_{d,i}}
 =\frac{(-\rho;\rho)_d}{(\rho;\rho)_d}.}
\end{equation}
At \(\rho=1/4\),
\begin{equation}\label{eq:stability-numeric}
 \Lambda_d(1/4)<\frac{(-1/4;1/4)_\infty}{(1/4;1/4)_\infty}
 =1.9692603536682694319\ldots.
\end{equation}
Consequently, if each sample is perturbed by at most \(\varepsilon\), the
extracted value is perturbed by less than \(1.969261\,\varepsilon\).
\end{proposition}

\begin{proof}
Put \(k=d-i\) in the absolute-value sum and use
\eqref{eq:single-extraction} without its alternating signs:
\begin{align*}
 \Lambda_d(\rho)
 &=\frac1{(\rho;\rho)_d}
 \sum_{k=0}^{d}\rho^{k(k+1)/2}\qbinom{d}{k}{\rho}\\
 &=\frac{(-\rho;\rho)_d}{(\rho;\rho)_d},
\end{align*}
where the second equality is the finite \(q\)-binomial theorem applied to
\((-\rho;\rho)_d\).  The ratio increases to the displayed convergent infinite-product limit when
\(\rho=1/4\).  The perturbation bound is
the triangle inequality.
\end{proof}

This bounded factor is a useful contrast with ordinary high-order Richardson
extrapolation at nearly coincident nodes, whose weights can be enormous.  The
geometric nodes \(1,1/4,1/16,\ldots\) remain well separated on the logarithmic
scale.

\section{Recovering every geometric coefficient and even derivative}

The same sample window determines all coefficients, not only the constant
one.  Apply to \eqref{eq:geometric-tail} the shift polynomial that kills every
mode except the \(s\)-th.

\begin{theorem}[Mode projector]\label{thm:mode-projector}
For \(0\le s\le d\) and \(N\ge m\),
\begin{equation}\label{eq:mode-projector}
 \boxed{\displaystyle
 C_s(x)=\rho^{-sN}
 \frac{\left[\displaystyle\prod_{\substack{0\le j\le d\\j\ne s}}
 (\E-\rho^j)q\right]_N}
 {\displaystyle\prod_{\substack{0\le j\le d\\j\ne s}}
 (\rho^s-\rho^j)},}
\end{equation}
where \(C_0(x)=\up(x)\).  The denominator has the closed form
\begin{equation}\label{eq:mode-denominator}
 \boxed{\displaystyle
 \prod_{j\ne s}(\rho^s-\rho^j)
 =(-1)^s\rho^{s(s-1)/2+s(d-s)}
 (\rho;\rho)_s(\rho;\rho)_{d-s}.}
\end{equation}
\end{theorem}

\begin{proof}
On the mode \(C_r\rho^{rn}\), the numerator operator acts by multiplication
with \(\prod_{j\ne s}(\rho^r-\rho^j)\).  This is zero unless \(r=s\).  For
\(r=s\), division by the displayed eigenvalue and by \(\rho^{sN}\) recovers
\(C_s\).

For \(j<s\), write
\(\rho^s-\rho^j=-\rho^j(1-\rho^{s-j})\); for \(j>s\), write
\(\rho^s-\rho^j=\rho^s(1-\rho^{j-s})\).  Multiplication gives
\eqref{eq:mode-denominator}.
\end{proof}

For a completely expanded \(q\)-binomial formula, split the numerator into
roots below and above \(s\).  The finite \(q\)-binomial theorem gives
\begin{align}\label{eq:mode-numerator-expanded}
 &\left[\prod_{j\ne s}(\E-\rho^j)q\right]_N\\
 &\quad=\sum_{a=0}^{s}\sum_{b=0}^{d-s}
 (-1)^{a+b}
 \rho^{a(a-1)/2+(s+1)b+b(b-1)/2}
 \qbinom{s}{a}{\rho}\qbinom{d-s}{b}{\rho}
 q_{N+d-a-b}.\nonumber
\end{align}
Equations \eqref{eq:mode-projector}, \eqref{eq:mode-denominator}, and
\eqref{eq:mode-numerator-expanded} are a closed-form inverse for the entire
geometric Vandermonde system.

Finally, \eqref{eq:Cj-def} converts geometric modes into the even derivative
jet.

\begin{corollary}[Even-jet recovery]\label{cor:jet-recovery}
For \(1\le s\le d\),
\begin{equation}\label{eq:jet-recovery}
 \boxed{\displaystyle
 \up^{(2s)}(x)=\frac{(-1)^s4^s}{a_s}\,C_s(x).}
\end{equation}
Thus \(d+1\) consecutive finite-spline values recover not only \(\up(x)\)
but every even derivative \(\up^{(2)},\ldots,\up^{(2d)}\) that can occur in
the terminating formula.
\end{corollary}

\section{The exact \texorpdfstring{$q$}{q}-difference recurrence}

Because \(q_n\) is a finite linear combination of the modes
\(1,\rho^n,\ldots,\rho^{dn}\), it satisfies a constant-coefficient recurrence.

\begin{theorem}[Annihilating recurrence]\label{thm:recurrence}
For every \(n\ge m\),
\begin{equation}\label{eq:operator-recurrence}
 \prod_{j=0}^{d}(\E-\rho^j)q_n=0.
\end{equation}
Equivalently,
\begin{equation}\label{eq:qbinomial-recurrence}
 \boxed{\displaystyle
 \sum_{k=0}^{d+1}(-1)^k\rho^{k(k-1)/2}
 \qbinom{d+1}{k}{\rho}\,
 q_{n+d+1-k}=0.}
\end{equation}
For odd \(m\), the recurrence already applies to the window beginning at
\(n=m-1\) under \cref{prop:odd-onset}.
\end{theorem}

\begin{proof}
Each factor \(\E-\rho^j\) annihilates its corresponding mode.  Expanding
\[
 \prod_{j=0}^{d}(\E-\rho^j)
 =\E^{d+1}(\E^{-1};\rho)_{d+1}
\]
by the finite \(q\)-binomial theorem gives
\eqref{eq:qbinomial-recurrence}.
\end{proof}

The recurrence provides a simple exact diagnostic.  If rationally computed
samples fail \eqref{eq:qbinomial-recurrence}, either the starting stage is too
early, the indexing convention is inconsistent, or an arithmetic error has
occurred.  In floating-point computation, the recurrence residual gives an
internal estimate of numerical contamination.

\section{Exact examples}

All fractions in this section are produced by the truncated-power formula
\eqref{eq:truncated-power}; no numerical approximation is involved.

\begin{example}[Stationary levels \(m=0,1\)]
At \(x=0\), one has \(m=d=0\), and
\[
 \up_n(0)=\up(0)=1\qquad(n\ge0).
\]
At \(x=\pm1/2\), one has \(m=1,d=0\), and
\[
 \up_n(\pm1/2)=\up(\pm1/2)=\frac12\qquad(n\ge1).
\]
Thus no decaying terms are present at levels zero and one.
\end{example}

\begin{example}[The point \(x=1/4\)]\label{ex:quarter}
Here \(m=2,d=1\).  The first guaranteed samples are
\begin{equation}\label{eq:quarter-samples}
 q_2=\frac{15}{16},
 \qquad
 q_3=\frac{179}{192}.
\end{equation}
Formula \eqref{eq:weights-d1} gives
\begin{equation}\label{eq:quarter-value}
 \up\!\left(\frac14\right)
 =\frac{-q_2+4q_3}{3}
 =\frac{67}{72}.
\end{equation}
The derivative formula gives \(\up''(1/4)=-8\), and \(a_1=1/18\).  Therefore
\begin{equation}\label{eq:quarter-law}
 \boxed{\displaystyle
 q_n=\frac{67}{72}+\frac19\,4^{-n},\qquad n\ge2.}
\end{equation}
The preceding value \(q_1=1\) does not satisfy this law, so the guaranteed
onset \(n=m=2\) is sharp in this example.
\end{example}

\begin{example}[Odd level: \(x=1/8\)]\label{ex:eighth}
Here \(m=3,d=1\).  The ordinary centered-cell theorem starts at \(n=3\), but
\cref{prop:odd-onset} moves the onset to \(n=2\).  Indeed,
\begin{equation}\label{eq:eighth-samples}
 q_2=1,
 \qquad q_3=\frac{383}{384},
 \qquad q_4=\frac{1531}{1536}.
\end{equation}
Using the first two gives
\begin{equation}\label{eq:eighth-value}
 \up\!\left(\frac18\right)
 =\frac{-q_2+4q_3}{3}
 =\frac{287}{288}.
\end{equation}
Since \(\up''(1/8)=-4\),
\begin{equation}\label{eq:eighth-law}
 \boxed{\displaystyle
 q_n=\frac{287}{288}+\frac1{18}\,4^{-n},\qquad n\ge2.}
\end{equation}
The value of \(q_4\) follows automatically from the same law.
\end{example}

\begin{example}[Two modes: \(x=1/16\)]\label{ex:sixteenth}
Now \(m=4,d=2\), so three samples are required:
\begin{equation}\label{eq:sixteenth-samples}
 q_4=\frac{24575}{24576},\qquad
 q_5=\frac{1965959}{1966080},\qquad
 q_6=\frac{94365509}{94371840}.
\end{equation}
The \(d=2\) extraction gives
\begin{equation}\label{eq:sixteenth-value}
 \begin{aligned}
 \up\!\left(\frac1{16}\right)
 &=\frac{q_4-20q_5+64q_6}{45}\\
 &=\frac{2073457}{2073600}.
 \end{aligned}
\end{equation}
The even derivatives are
\begin{equation}\label{eq:sixteenth-derivatives}
 \up''\!\left(\frac1{16}\right)=-\frac59,
 \qquad
 \up^{(4)}\!\left(\frac1{16}\right)=-1024.
\end{equation}
Using \(a_1=1/18\) and \(a_2=31/16200\) gives the complete exact sequence
\begin{equation}\label{eq:sixteenth-law}
 \boxed{\displaystyle
 q_n=\frac{2073457}{2073600}
 +\frac5{648}\,4^{-n}
 -\frac{248}{2025}\,16^{-n},
 \qquad n\ge4.}
\end{equation}
This example displays the two distinct geometric ratios \(1/4\) and \(1/16\)
and illustrates why ordinary one-step Richardson extrapolation is not enough
at higher dyadic levels.
\end{example}

\begin{example}[Reflection]
Every finite product and the limiting density are even.  Hence
\[
 \up_n(-x)=\up_n(x),\qquad \up(-x)=\up(x),
\]
and all formulas above apply unchanged after replacing \(x\) by \(-x\).
At \(x=\pm1\), every finite approximant vanishes because its support is
strictly contained in \([-1,1]\), and the smooth limiting function also
vanishes.  The endpoint sequence is identically zero.
\end{example}

\section{An exact extraction algorithm}

The mathematical procedure separates cleanly into two layers.  The first
layer computes rational finite-spline values by any exact method; the second
layer applies a short, universal \(q\)-Pochhammer filter.

\begin{algorithm}[Dyadic limit extraction]\label{alg:extraction}
Given a dyadic rational \(x\in[-1,1]\):
\begin{enumerate}
 \item If \(x=\pm1\), return \(0\).
 \item Reduce \(x\) to lowest terms and let \(m\) be the base-two logarithm of
 its denominator.  Put \(d=\floor{m/2}\).
 \item Choose a starting stage \(N\ge m\).  The canonical choice is \(N=m\).
 For odd \(m\ge3\), \(N=m-1\) is also valid by \cref{prop:odd-onset}.
 \item Compute the \(d+1\) exact rational values
 \[
  q_N,q_{N+1},\ldots,q_{N+d}.
 \]
 Formula \eqref{eq:truncated-power} is a direct reference implementation.
 \item Compute the weights from either \eqref{eq:forward-extraction} or
 \eqref{eq:integer-weight}, and return
 \[
  \sum_{i=0}^{d}w_{d,i}q_{N+i}.
 \]
\end{enumerate}
The returned fraction is exactly \(\up(x)\).
\end{algorithm}

\subsection{Pseudocode}

\begin{lstlisting}[style=pythonstyle]
def extract_up_at_dyadic(x):
    # x is an exact reduced Fraction with power-of-two denominator.
    if abs(x) == 1:
        return Fraction(0)

    m = log2(x.denominator)       # exact integer
    d = m // 2
    rho = Fraction(1, 4)
    N = m                         # unconditional starting stage

    samples = [finite_up(N + i, x) for i in range(d + 1)]
    weights = []
    for i in range(d + 1):
        T = (d - i) * (d - i + 1) // 2
        w = (-1)**(d - i) * rho**T
        w /= q_pochhammer(rho, i) * q_pochhammer(rho, d - i)
        weights.append(w)

    return sum(w * q for w, q in zip(weights, samples))
\end{lstlisting}

\subsection{Cost and validation}

The direct truncated-power evaluator has \(2^{n+1}\) potential terms at stage
\(n\), though only the active prefix is needed.  It is ideal for transparent
verification at modest levels.  Faster production evaluators can exploit
binary block recurrences, local spline propagation, or precomputed dyadic
rows; none of those choices changes the extraction filter.

Once the samples are available, computing the weights and their dot product
requires only \(O(d)\) rational operations if successive products
\((\rho;\rho)_i\) are cached.  Since \(d=\floor{m/2}\), this second layer is
small even for fairly deep dyadic levels.

Two exact checks are especially useful:
\begin{enumerate}
 \item repeat the extraction with the shifted window
 \(q_{N+1},\ldots,q_{N+d+1}\); the resulting rational must be identical;
 \item evaluate the recurrence \eqref{eq:qbinomial-recurrence}; its residual
 must be exactly zero.
\end{enumerate}
If an upper bound \(D\ge d\) is used instead of the exact \(d\), the
\((D+1)\)-sample filter still returns the same constant mode, because it
annihilates all modes through \(\rho^{Dn}\).  Underestimating \(d\), by
contrast, generally leaves an uncancelled geometric term.

\subsection{Approximate samples and rational reconstruction}

When \(q_n\) is evaluated numerically rather than rationally, apply the same
weights at high precision.  By \cref{prop:stability}, the extraction amplifies
uniform absolute error by less than a factor of \(1.97\).  If a denominator
bound is known from independent arithmetic theory, one may then use continued
fractions or rational reconstruction on the high-precision result.  Exact
input remains preferable: formula \eqref{eq:truncated-power} and the extraction
weights involve rational arithmetic only.

\section{Companion implementation and verification}

The archive accompanying this article contains
\texttt{up\_dyadic\_extraction.py}.  It uses only the Python standard library
and \texttt{fractions.Fraction}.  The program implements:
\begin{itemize}
 \item the exact truncated-power evaluator \eqref{eq:truncated-power};
 \item finite \(q\)-Pochhammer symbols, Gaussian coefficients, and extraction
 weights;
 \item the Bell--Bernoulli recurrence \eqref{eq:a-recurrence};
 \item exact dyadic evaluation of \(\up\) by \eqref{eq:single-extraction};
 \item exact derivative evaluation by \eqref{eq:derivative-formula};
 \item exhaustive theorem checks through a user-selected binary level.
\end{itemize}

For example,
\begin{lstlisting}[style=pythonstyle]
python up_dyadic_extraction.py --examples
python up_dyadic_extraction.py --x 1/16
python up_dyadic_extraction.py --verify-level 7
\end{lstlisting}
The bundled verification run checked \(2376\) exact identities at all \(255\)
interior dyadic points of exact level at most seven.  This includes sliding-window
invariance of the \(q\)-Pochhammer extraction and the differential formula
\eqref{eq:main-differential} at several consecutive stages.  Every comparison
was an equality of rational numbers; no tolerance or floating-point arithmetic
was used.

The computation is evidence that the formulas have been transcribed and
indexed correctly.  The proofs in the preceding sections are independent of
the experiment.

\section{Conceptual consequences and extensions}

\subsection{A second proof of dyadic rationality}

Classical arithmetic formulas already imply that \(\up(x)\) is rational at
dyadic \(x\).  The present construction gives a short independent route.
Every \(q_{N+i}\) is rational by \cref{cor:rational-samples}; every weight in
\eqref{eq:forward-extraction} is rational; hence their finite sum
\(\up(x)\) is rational.  The argument is constructive and returns the value
itself, not merely a rationality certificate.

\subsection{Why the number of modes is half the binary level}

The local polynomial degree is at most \(m\), but the omitted tail is symmetric.
Its moment-generating function is even, so the deconvolution operator contains
only even powers \(D^{2j}\).  Therefore only
\[
 0,2,4,\ldots,2\floor{m/2}
\]
can contribute at the point.  This is why the geometric degree is
\(d=\floor{m/2}\), rather than \(m\), and why adding one binary denominator
bit increases the sample count only every other time.

\subsection{Universality of the \texorpdfstring{$q$}{q}-filter}

The extraction theorem is not specific to the up-function once the finite
geometric law has been established.  If any sequence has the form
\begin{equation}\label{eq:universal-geometric}
 s_n=L+\sum_{j=1}^{d}A_j\rho^{jn},
\end{equation}
then exactly the same weights recover \(L\) from any \(d+1\) consecutive
samples.  The up-function supplies an unusually rigid mechanism that proves
\eqref{eq:universal-geometric} exactly: a self-similar omitted tail, a local
polynomial, and Prouhet cancellation.

\subsection{Beyond point values}

Corollary \ref{cor:derivative-tail} already extends the method to derivative
samples.  More generally, any linear functional depending only on the local
polynomial jet inherits a finite geometric law by applying that functional to
\eqref{eq:polynomial-inverse}.  This includes finite combinations of values
and derivatives inside a fixed centered cell.  The same algebra can therefore
serve as an exact extrapolation layer for local approximation schemes built
from the finite dyadic box splines.

\subsection{Relation to ordinary Richardson extrapolation}

Ordinary Richardson extrapolation assumes an expansion
\(s_n=L+c_1h_n+c_2h_n^2+\cdots\) and cancels finitely many leading terms.
Here \(h_n=4^{-n}\), but two stronger statements hold:
\begin{enumerate}
 \item the admissible exponents \(1,2,\ldots,d\) are known exactly from the
 binary level of \(x\);
 \item after exponent \(d\), there is no remainder at all.
\end{enumerate}
Thus \eqref{eq:single-extraction} is simultaneously a geometric Richardson
formula, a Lagrange-evaluation formula at zero, a spectral projector for the
shift operator, and an exact inverse-Vandermonde row.

\section{Conclusion}

For a fixed dyadic rational \(x\), the finite-sinc approximants do not merely
approach \(\up(x)\) with an asymptotic power series.  After finitely many
initial stages, their values lie exactly in the finite-dimensional span
\[
 1,4^{-n},16^{-n},\ldots,4^{-dn},
 \qquad d=\floor{\ell(x)/2}.
\]
The mechanism is structural: the binary denominator aligns \(x\) with a
finite-spline cell, complete Thue--Morse blocks collapse its local degree, and
the self-similar omitted tail acts through an even moment operator whose
inverse terminates on that polynomial.

The resulting limit extraction is explicit, rational, stable, and compact:
\[
 \up(x)=\frac{1}{(1/4;1/4)_d}
 \sum_{k=0}^{d}(-1)^k4^{-k(k+1)/2}
 \qbinom{d}{k}{1/4}\,\up_{N+d-k}(x).
\]
The same sample window recovers every geometric coefficient and every relevant
even derivative, while the finite \(q\)-binomial recurrence supplies an exact
consistency check.  In this setting, \(q\)-Pochhammer extrapolation is not a
numerical acceleration heuristic; it is an exact arithmetic identity.

\appendix

\section{Derivation of the integer-product weights}

For \(\rho=1/4\),
\begin{equation}\label{eq:qpoch-P}
 (\rho;\rho)_s
 =\prod_{r=1}^{s}(1-4^{-r})
 =4^{-s(s+1)/2}P_s.
\end{equation}
Substitute \eqref{eq:qpoch-P} into \eqref{eq:forward-extraction}:
\begin{align*}
 w_{d,i}
 &=(-1)^{d-i}
 \frac{4^{-(d-i)(d-i+1)/2}}
 {4^{-i(i+1)/2}P_i\,
  4^{-(d-i)(d-i+1)/2}P_{d-i}}\\
 &=(-1)^{d-i}\frac{4^{i(i+1)/2}}{P_iP_{d-i}}.
\end{align*}
This proves \eqref{eq:integer-weight}.  It also explains the conspicuous
powers
\[
 1,4,64,4096,1048576,\ldots=4^{i(i+1)/2}
\]
in the rightmost numerators of the low-order formulas.

\section{A compact proof of the finite \texorpdfstring{$q$}{q}-binomial identity}

For completeness, define
\[
 (z;\rho)_d=\prod_{j=0}^{d-1}(1-z\rho^j).
\]
Induction on \(d\), using the Gaussian Pascal recurrence, gives
\begin{equation}\label{eq:finite-qbinomial-appendix}
 (z;\rho)_d
 =\sum_{k=0}^{d}(-1)^k\rho^{k(k-1)/2}
 \qbinom{d}{k}{\rho}z^k.
\end{equation}
Taking \(z=\rho\E^{-1}\) and multiplying by \(\E^d\) yields the extraction
operator expansion.  Taking \(z=\E^{-1}\) and multiplying by
\(\E^{d+1}\) yields the recurrence.  Taking \(z=-\rho\) yields the absolute
weight sum in \cref{prop:stability}.

\section{Reference formulas in one place}

For convenient implementation, the main identities are collected here.
Let \(x\in(-1,1)\) have exact level \(m\), let
\(d=\floor{m/2}\), \(\rho=1/4\), and choose \(N\ge m\).

\paragraph{Finite spline.}
\[
 q_n=\frac{2^{n(n+1)/2}}{n!}
 \sum_{k=0}^{2^{n+1}-1}\tau_k
 \pos{x+1-2^{-n-1}-k2^{-n}}^n.
\]

\paragraph{Exact geometric law.}
\[
 q_n=\sum_{j=0}^{d}(-1)^ja_j4^{-j(n+1)}\up^{(2j)}(x)
 =\up(x)+\sum_{j=1}^{d}C_j(x)\rho^{jn}.
\]

\paragraph{Bell--Bernoulli coefficients.}
\[
 \alpha_j=\frac{\abs{B_{2j}}}{2j(2j)!(1-2^{-2j})},
 \qquad
 ja_j=\sum_{r=1}^{j}r\alpha_ra_{j-r},
 \qquad a_0=1.
\]

\paragraph{Limit extraction.}
\[
 \up(x)=\frac1{(\rho;\rho)_d}
 \sum_{k=0}^{d}(-1)^k\rho^{k(k+1)/2}
 \qbinom{d}{k}{\rho}q_{N+d-k}.
\]

\paragraph{Mode extraction.}
\[
 C_s=\rho^{-sN}
 \frac{\left[\prod_{j\ne s}(\E-\rho^j)q\right]_N}
 {(-1)^s\rho^{s(s-1)/2+s(d-s)}
  (\rho;\rho)_s(\rho;\rho)_{d-s}}.
\]

\paragraph{Derivative recovery.}
\[
 \up^{(2s)}(x)=\frac{(-1)^s4^s}{a_s}C_s.
\]

\paragraph{Recurrence.}
\[
 \sum_{k=0}^{d+1}(-1)^k\rho^{k(k-1)/2}
 \qbinom{d+1}{k}{\rho}q_{n+d+1-k}=0.
\]

\begin{thebibliography}{99}

\bibitem{ProveIt}
V. Reshetnikov,
\emph{ProveIt: Analysis/FabiusFunction}, especially the finite dyadic sinc
product, exact box-spline, Bell--Bernoulli deconvolution, and geometric
Richardson developments,
\url{https://github.com/VladimirReshetnikov/ProveIt/tree/main/Analysis/FabiusFunction},
repository consulted 2 September 2026.

\bibitem{AriasArithmetic}
J. Arias de Reyna,
\emph{Arithmetic of the Fabius function},
2017,
\url{https://arxiv.org/abs/1702.06487}.

\bibitem{AriasCompact}
J. Arias de Reyna,
\emph{An infinitely differentiable function with compact support: Definition
and properties},
English translation of the 1982 article,
2017,
\url{https://arxiv.org/abs/1702.05442}.

\bibitem{Haugland}
J. K. Haugland,
\emph{Evaluating the Fabius function},
2016,
\url{https://arxiv.org/abs/1609.07999}.

\bibitem{DLMFq}
NIST Digital Library of Mathematical Functions,
\emph{Chapter 17: \(q\)-Hypergeometric and Related Functions}, especially
Section 17.2 on \(q\)-Pochhammer symbols, Gaussian coefficients, and the
finite \(q\)-binomial theorem,
\url{https://dlmf.nist.gov/17.2}.

\end{thebibliography}

\end{document}
